\documentclass[aspectratio=169,11pt]{beamer}
% Lecture 02 — Int, division, mutable vars, floats, Long/Short/Byte
% Course: Introduction to Programming with Kotlin

% Shared Beamer preamble for Introduction to Programming with Kotlin
% Include with: % Shared Beamer preamble for Introduction to Programming with Kotlin
% Include with: % Shared Beamer preamble for Introduction to Programming with Kotlin
% Include with: \input{preamble}

\usetheme{Madrid}
\usecolortheme{default}

% Clean, professional palette (deep navy + muted teal accent)
\definecolor{LectureNavy}{HTML}{1B3A4B}
\definecolor{LectureAccent}{HTML}{2A9D8F}
\definecolor{LectureGray}{HTML}{4A5568}
\definecolor{CodeBg}{HTML}{F7FAFC}
\definecolor{AlertBg}{HTML}{FFF5F5}
\definecolor{TryBg}{HTML}{F0FFF4}
\definecolor{QuizBg}{HTML}{EBF8FF}
\definecolor{AnswerKeyBg}{HTML}{FFFFF0}
\definecolor{AnswerGold}{HTML}{B7791F}
\definecolor{KeywordBlue}{HTML}{2B6CB0}
\definecolor{StringGreen}{HTML}{276749}
\definecolor{CommentGray}{HTML}{718096}

\setbeamercolor{structure}{fg=LectureNavy}
\setbeamercolor{palette primary}{bg=LectureNavy,fg=white}
\setbeamercolor{palette secondary}{bg=LectureAccent,fg=white}
\setbeamercolor{palette tertiary}{bg=LectureGray,fg=white}
% Madrid puts title/subtitle in a coloured box (inherits palette primary bg).
% Title text must be light on that dark box — not navy-on-navy.
\setbeamercolor{title}{fg=white,bg=LectureNavy}
\setbeamercolor{subtitle}{fg=white,bg=LectureNavy}
\setbeamercolor{author}{fg=LectureNavy}
\setbeamercolor{date}{fg=LectureGray}
\setbeamercolor{institute}{fg=LectureGray}
\setbeamercolor{frametitle}{bg=LectureNavy,fg=white}
\setbeamercolor{block title}{bg=LectureNavy,fg=white}
\setbeamercolor{block body}{bg=CodeBg,fg=black}
\setbeamercolor{block title alerted}{bg=LectureAccent,fg=white}
\setbeamercolor{block body alerted}{bg=TryBg,fg=black}
\setbeamercolor{block title example}{bg=LectureGray,fg=white}
\setbeamercolor{block body example}{bg=AlertBg,fg=black}

\setbeamertemplate{navigation symbols}{}
\setbeamertemplate{itemize items}[circle]
\setbeamertemplate{enumerate items}[default]

\usepackage[T1]{fontenc}
\usepackage[utf8]{inputenc}
\usepackage{lmodern}
\usepackage{booktabs}
\usepackage{tabularx}
\usepackage{graphicx}
\usepackage{hyperref}
\usepackage{listings}
\usepackage{xcolor}
\usepackage{tikz}
\usetikzlibrary{positioning,arrows.meta}

% listings: Kotlin without shell-escape (no built-in Kotlin language)
\lstdefinelanguage{Kotlin}{
  morekeywords={
    abstract,actual,annotation,as,break,by,catch,class,companion,const,
    constructor,continue,crossinline,data,do,dynamic,else,enum,expect,
    external,false,field,file,final,finally,for,fun,get,if,import,in,
    infix,init,inline,inner,interface,internal,is,lateinit,noinline,null,
    object,open,operator,out,override,package,private,protected,public,
    reified,return,sealed,set,super,suspend,tailrec,this,throw,true,try,
    typealias,typeof,val,var,when,where,while
  },
  sensitive=true,
  morecomment=[l]{//},
  morecomment=[s]{/*}{*/},
  morestring=[b]",
  morestring=[b]',
  morestring=[s]{"""}{"""},
}

\lstdefinestyle{kotlinstyle}{
  language=Kotlin,
  basicstyle=\ttfamily\small,
  keywordstyle=\color{KeywordBlue}\bfseries,
  commentstyle=\color{CommentGray}\itshape,
  stringstyle=\color{StringGreen},
  numbers=left,
  numberstyle=\tiny\color{CommentGray},
  numbersep=6pt,
  frame=single,
  rulecolor=\color{LectureGray!40},
  backgroundcolor=\color{CodeBg},
  breaklines=true,
  showstringspaces=false,
  tabsize=2,
  captionpos=b,
  morekeywords={readln,readLine,println,print,listOf,mutableListOf,toInt,toDouble,toLong,toShort,toByte,toChar,toString,split,size},
}
\lstset{style=kotlinstyle}

\newcommand{\kotlincode}[1]{\texttt{\small #1}}
\newcommand{\trythis}[1]{%
  \begin{alertblock}{Try this}
    #1
  \end{alertblock}%
}
\newcommand{\commonmistake}[1]{%
  \begin{exampleblock}{Common mistake}
    #1
  \end{exampleblock}%
}
% Formative in-class checks: question slide, then a dedicated Answer slide
\newcommand{\quizpause}{%
  \par\vspace{0.5em}
  {\small\textit{Pause --- answer (clicker / raise hand / neighbour) before the next slide.}}%
}
\newcommand{\quizbox}[1]{%
  \setbeamercolor{block title}{bg=KeywordBlue,fg=white}%
  \setbeamercolor{block body}{bg=QuizBg,fg=black}%
  \begin{block}{Quiz}
    #1
  \end{block}%
  \setbeamercolor{block title}{bg=LectureNavy,fg=white}%
  \setbeamercolor{block body}{bg=CodeBg,fg=black}%
}
\newcommand{\answerbox}[1]{%
  \setbeamercolor{block title}{bg=AnswerGold,fg=white}%
  \setbeamercolor{block body}{bg=AnswerKeyBg,fg=black}%
  \begin{block}{Answer}
    #1
  \end{block}%
  \setbeamercolor{block title}{bg=LectureNavy,fg=white}%
  \setbeamercolor{block body}{bg=CodeBg,fg=black}%
}

% Empty author/date placeholders for the lecturer to fill
\author{}
\date{}


\usetheme{Madrid}
\usecolortheme{default}

% Clean, professional palette (deep navy + muted teal accent)
\definecolor{LectureNavy}{HTML}{1B3A4B}
\definecolor{LectureAccent}{HTML}{2A9D8F}
\definecolor{LectureGray}{HTML}{4A5568}
\definecolor{CodeBg}{HTML}{F7FAFC}
\definecolor{AlertBg}{HTML}{FFF5F5}
\definecolor{TryBg}{HTML}{F0FFF4}
\definecolor{QuizBg}{HTML}{EBF8FF}
\definecolor{AnswerKeyBg}{HTML}{FFFFF0}
\definecolor{AnswerGold}{HTML}{B7791F}
\definecolor{KeywordBlue}{HTML}{2B6CB0}
\definecolor{StringGreen}{HTML}{276749}
\definecolor{CommentGray}{HTML}{718096}

\setbeamercolor{structure}{fg=LectureNavy}
\setbeamercolor{palette primary}{bg=LectureNavy,fg=white}
\setbeamercolor{palette secondary}{bg=LectureAccent,fg=white}
\setbeamercolor{palette tertiary}{bg=LectureGray,fg=white}
% Madrid puts title/subtitle in a coloured box (inherits palette primary bg).
% Title text must be light on that dark box — not navy-on-navy.
\setbeamercolor{title}{fg=white,bg=LectureNavy}
\setbeamercolor{subtitle}{fg=white,bg=LectureNavy}
\setbeamercolor{author}{fg=LectureNavy}
\setbeamercolor{date}{fg=LectureGray}
\setbeamercolor{institute}{fg=LectureGray}
\setbeamercolor{frametitle}{bg=LectureNavy,fg=white}
\setbeamercolor{block title}{bg=LectureNavy,fg=white}
\setbeamercolor{block body}{bg=CodeBg,fg=black}
\setbeamercolor{block title alerted}{bg=LectureAccent,fg=white}
\setbeamercolor{block body alerted}{bg=TryBg,fg=black}
\setbeamercolor{block title example}{bg=LectureGray,fg=white}
\setbeamercolor{block body example}{bg=AlertBg,fg=black}

\setbeamertemplate{navigation symbols}{}
\setbeamertemplate{itemize items}[circle]
\setbeamertemplate{enumerate items}[default]

\usepackage[T1]{fontenc}
\usepackage[utf8]{inputenc}
\usepackage{lmodern}
\usepackage{booktabs}
\usepackage{tabularx}
\usepackage{graphicx}
\usepackage{hyperref}
\usepackage{listings}
\usepackage{xcolor}
\usepackage{tikz}
\usetikzlibrary{positioning,arrows.meta}

% listings: Kotlin without shell-escape (no built-in Kotlin language)
\lstdefinelanguage{Kotlin}{
  morekeywords={
    abstract,actual,annotation,as,break,by,catch,class,companion,const,
    constructor,continue,crossinline,data,do,dynamic,else,enum,expect,
    external,false,field,file,final,finally,for,fun,get,if,import,in,
    infix,init,inline,inner,interface,internal,is,lateinit,noinline,null,
    object,open,operator,out,override,package,private,protected,public,
    reified,return,sealed,set,super,suspend,tailrec,this,throw,true,try,
    typealias,typeof,val,var,when,where,while
  },
  sensitive=true,
  morecomment=[l]{//},
  morecomment=[s]{/*}{*/},
  morestring=[b]",
  morestring=[b]',
  morestring=[s]{"""}{"""},
}

\lstdefinestyle{kotlinstyle}{
  language=Kotlin,
  basicstyle=\ttfamily\small,
  keywordstyle=\color{KeywordBlue}\bfseries,
  commentstyle=\color{CommentGray}\itshape,
  stringstyle=\color{StringGreen},
  numbers=left,
  numberstyle=\tiny\color{CommentGray},
  numbersep=6pt,
  frame=single,
  rulecolor=\color{LectureGray!40},
  backgroundcolor=\color{CodeBg},
  breaklines=true,
  showstringspaces=false,
  tabsize=2,
  captionpos=b,
  morekeywords={readln,readLine,println,print,listOf,mutableListOf,toInt,toDouble,toLong,toShort,toByte,toChar,toString,split,size},
}
\lstset{style=kotlinstyle}

\newcommand{\kotlincode}[1]{\texttt{\small #1}}
\newcommand{\trythis}[1]{%
  \begin{alertblock}{Try this}
    #1
  \end{alertblock}%
}
\newcommand{\commonmistake}[1]{%
  \begin{exampleblock}{Common mistake}
    #1
  \end{exampleblock}%
}
% Formative in-class checks: question slide, then a dedicated Answer slide
\newcommand{\quizpause}{%
  \par\vspace{0.5em}
  {\small\textit{Pause --- answer (clicker / raise hand / neighbour) before the next slide.}}%
}
\newcommand{\quizbox}[1]{%
  \setbeamercolor{block title}{bg=KeywordBlue,fg=white}%
  \setbeamercolor{block body}{bg=QuizBg,fg=black}%
  \begin{block}{Quiz}
    #1
  \end{block}%
  \setbeamercolor{block title}{bg=LectureNavy,fg=white}%
  \setbeamercolor{block body}{bg=CodeBg,fg=black}%
}
\newcommand{\answerbox}[1]{%
  \setbeamercolor{block title}{bg=AnswerGold,fg=white}%
  \setbeamercolor{block body}{bg=AnswerKeyBg,fg=black}%
  \begin{block}{Answer}
    #1
  \end{block}%
  \setbeamercolor{block title}{bg=LectureNavy,fg=white}%
  \setbeamercolor{block body}{bg=CodeBg,fg=black}%
}

% Empty author/date placeholders for the lecturer to fill
\author{}
\date{}


\usetheme{Madrid}
\usecolortheme{default}

% Clean, professional palette (deep navy + muted teal accent)
\definecolor{LectureNavy}{HTML}{1B3A4B}
\definecolor{LectureAccent}{HTML}{2A9D8F}
\definecolor{LectureGray}{HTML}{4A5568}
\definecolor{CodeBg}{HTML}{F7FAFC}
\definecolor{AlertBg}{HTML}{FFF5F5}
\definecolor{TryBg}{HTML}{F0FFF4}
\definecolor{QuizBg}{HTML}{EBF8FF}
\definecolor{AnswerKeyBg}{HTML}{FFFFF0}
\definecolor{AnswerGold}{HTML}{B7791F}
\definecolor{KeywordBlue}{HTML}{2B6CB0}
\definecolor{StringGreen}{HTML}{276749}
\definecolor{CommentGray}{HTML}{718096}

\setbeamercolor{structure}{fg=LectureNavy}
\setbeamercolor{palette primary}{bg=LectureNavy,fg=white}
\setbeamercolor{palette secondary}{bg=LectureAccent,fg=white}
\setbeamercolor{palette tertiary}{bg=LectureGray,fg=white}
% Madrid puts title/subtitle in a coloured box (inherits palette primary bg).
% Title text must be light on that dark box — not navy-on-navy.
\setbeamercolor{title}{fg=white,bg=LectureNavy}
\setbeamercolor{subtitle}{fg=white,bg=LectureNavy}
\setbeamercolor{author}{fg=LectureNavy}
\setbeamercolor{date}{fg=LectureGray}
\setbeamercolor{institute}{fg=LectureGray}
\setbeamercolor{frametitle}{bg=LectureNavy,fg=white}
\setbeamercolor{block title}{bg=LectureNavy,fg=white}
\setbeamercolor{block body}{bg=CodeBg,fg=black}
\setbeamercolor{block title alerted}{bg=LectureAccent,fg=white}
\setbeamercolor{block body alerted}{bg=TryBg,fg=black}
\setbeamercolor{block title example}{bg=LectureGray,fg=white}
\setbeamercolor{block body example}{bg=AlertBg,fg=black}

\setbeamertemplate{navigation symbols}{}
\setbeamertemplate{itemize items}[circle]
\setbeamertemplate{enumerate items}[default]

\usepackage[T1]{fontenc}
\usepackage[utf8]{inputenc}
\usepackage{lmodern}
\usepackage{booktabs}
\usepackage{tabularx}
\usepackage{graphicx}
\usepackage{hyperref}
\usepackage{listings}
\usepackage{xcolor}
\usepackage{tikz}
\usetikzlibrary{positioning,arrows.meta}

% listings: Kotlin without shell-escape (no built-in Kotlin language)
\lstdefinelanguage{Kotlin}{
  morekeywords={
    abstract,actual,annotation,as,break,by,catch,class,companion,const,
    constructor,continue,crossinline,data,do,dynamic,else,enum,expect,
    external,false,field,file,final,finally,for,fun,get,if,import,in,
    infix,init,inline,inner,interface,internal,is,lateinit,noinline,null,
    object,open,operator,out,override,package,private,protected,public,
    reified,return,sealed,set,super,suspend,tailrec,this,throw,true,try,
    typealias,typeof,val,var,when,where,while
  },
  sensitive=true,
  morecomment=[l]{//},
  morecomment=[s]{/*}{*/},
  morestring=[b]",
  morestring=[b]',
  morestring=[s]{"""}{"""},
}

\lstdefinestyle{kotlinstyle}{
  language=Kotlin,
  basicstyle=\ttfamily\small,
  keywordstyle=\color{KeywordBlue}\bfseries,
  commentstyle=\color{CommentGray}\itshape,
  stringstyle=\color{StringGreen},
  numbers=left,
  numberstyle=\tiny\color{CommentGray},
  numbersep=6pt,
  frame=single,
  rulecolor=\color{LectureGray!40},
  backgroundcolor=\color{CodeBg},
  breaklines=true,
  showstringspaces=false,
  tabsize=2,
  captionpos=b,
  morekeywords={readln,readLine,println,print,listOf,mutableListOf,toInt,toDouble,toLong,toShort,toByte,toChar,toString,split,size},
}
\lstset{style=kotlinstyle}

\newcommand{\kotlincode}[1]{\texttt{\small #1}}
\newcommand{\trythis}[1]{%
  \begin{alertblock}{Try this}
    #1
  \end{alertblock}%
}
\newcommand{\commonmistake}[1]{%
  \begin{exampleblock}{Common mistake}
    #1
  \end{exampleblock}%
}
% Formative in-class checks: question slide, then a dedicated Answer slide
\newcommand{\quizpause}{%
  \par\vspace{0.5em}
  {\small\textit{Pause --- answer (clicker / raise hand / neighbour) before the next slide.}}%
}
\newcommand{\quizbox}[1]{%
  \setbeamercolor{block title}{bg=KeywordBlue,fg=white}%
  \setbeamercolor{block body}{bg=QuizBg,fg=black}%
  \begin{block}{Quiz}
    #1
  \end{block}%
  \setbeamercolor{block title}{bg=LectureNavy,fg=white}%
  \setbeamercolor{block body}{bg=CodeBg,fg=black}%
}
\newcommand{\answerbox}[1]{%
  \setbeamercolor{block title}{bg=AnswerGold,fg=white}%
  \setbeamercolor{block body}{bg=AnswerKeyBg,fg=black}%
  \begin{block}{Answer}
    #1
  \end{block}%
  \setbeamercolor{block title}{bg=LectureNavy,fg=white}%
  \setbeamercolor{block body}{bg=CodeBg,fg=black}%
}

% Empty author/date placeholders for the lecturer to fill
\author{}
\date{}


\title{Lecture 02: Numbers and Mutable Variables}
\subtitle{Introduction to Programming with Kotlin}

\begin{document}

%----------------------------------------------------------------------
    \begin{frame}
        \titlepage
    \end{frame}

%----------------------------------------------------------------------
    \begin{frame}{Agenda}
        \begin{enumerate}
            \item The integer type \kotlincode{Int}
            \item Integer division and remainder
            \item Mutable variables with \kotlincode{var}
            \item Floating-point numbers
            \item \kotlincode{Long}, \kotlincode{Short}, and \kotlincode{Byte}
        \end{enumerate}
    \end{frame}

%======================================================================
    \section{Integer type Int}
%======================================================================

    \begin{frame}{What is \kotlincode{Int}?}
        \kotlincode{Int} is Kotlin's default 32-bit signed integer type.
        \vspace{0.6em}
        \begin{itemize}
            \item Range: about $-2\cdot 10^9$ to $2\cdot 10^9$
              (exactly $-2^{31}$ \ldots\ $2^{31}-1$)
            \item Whole numbers only: $0$, $-3$, $42$, $1000$
            \item Integer literals like \kotlincode{42} have type \kotlincode{Int} by default
        \end{itemize}
        \vspace{0.5em}
        \begin{block}{When you need \kotlincode{Int}}
            Counts, ages, indices, years, loop counters, simple arithmetic on wholes.
        \end{block}
    \end{frame}

    \begin{frame}[fragile]{Declaring and using \kotlincode{Int}}
        \begin{lstlisting}
fun main() {
    val a: Int = 10
    val b = 3          // inferred as Int
    println(a + b)     // 13
    println(a - b)     // 7
    println(a * b)     // 30
}
        \end{lstlisting}
        \vspace{0.3em}
        Arithmetic operators: \kotlincode{+}, \kotlincode{-}, \kotlincode{*}, \kotlincode{/}, \kotlincode{\%}
    \end{frame}

    \begin{frame}[fragile]{Reading an \kotlincode{Int}}
        \begin{lstlisting}
fun main() {
    print("Enter an integer: ")
    val n = readln().toInt()
    println("You entered $n")
    println("Double is ${n * 2}")
}
        \end{lstlisting}
        \vspace{0.3em}
        \commonmistake{Forgetting \kotlincode{toInt()}.
        \kotlincode{readln()} returns \kotlincode{String}; \kotlincode{"10" + "20"} is
        \kotlincode{"1020"}, not $30$.}
    \end{frame}

    %--- Quiz checkpoint 1 ---
    \begin{frame}{Quiz 1 --- Int basics}
        \quizbox{What is the type of the literal \kotlincode{100} in Kotlin?}
        \begin{enumerate}
            \item \kotlincode{Long}
            \item \kotlincode{Int}
            \item \kotlincode{Double}
            \item \kotlincode{String}
        \end{enumerate}
        \quizpause
    \end{frame}

    \begin{frame}{Answer --- Quiz 1}
        \answerbox{\textbf{2.} \kotlincode{Int}}
        \vspace{0.4em}
        Unsuffixed integer literals are \kotlincode{Int}.
        Write \kotlincode{100L} for \kotlincode{Long}, or \kotlincode{100.0} for \kotlincode{Double}.
    \end{frame}

%======================================================================
    \section{Integer division}
%======================================================================

    \begin{frame}{Integer division}
        When both operands are integers, \kotlincode{/} performs \textbf{integer division}:
        the fractional part is discarded (truncated toward zero).
        \vspace{0.6em}
        \begin{itemize}
            \item \kotlincode{7 / 2} $\rightarrow$ \kotlincode{3} (not $3.5$)
            \item \kotlincode{(-7) / 2} $\rightarrow$ \kotlincode{-3}
            \item Remainder: \kotlincode{7 \% 2} $\rightarrow$ \kotlincode{1}
        \end{itemize}
    \end{frame}

    \begin{frame}[fragile]{Division and remainder}
        \begin{lstlisting}
fun main() {
    val total = 17
    val group = 5
    println(total / group)  // 3
    println(total % group)  // 2
    // 3 groups of 5, with 2 left over
}
        \end{lstlisting}
        \vspace{0.3em}
        Check: $(total / group) * group + (total \% group) = total$
        (for non-negative values).
    \end{frame}

    \begin{frame}[fragile]{Getting a fractional result}
        Convert at least one operand to floating-point:
        \vspace{0.3em}
        \begin{lstlisting}
fun main() {
    val a = 7
    val b = 2
    println(a / b)             // 3  (Int)
    println(a.toDouble() / b)  // 3.5
    println(7.0 / 2)           // 3.5
}
        \end{lstlisting}
        \trythis{Read two integers and print both the integer quotient and the
        exact ratio as a decimal.}
    \end{frame}

    %--- Quiz checkpoint 2 ---
    \begin{frame}{Quiz 2 --- Integer division}
        \quizbox{What does \kotlincode{println(20 / 3)} print?}
        \begin{enumerate}
            \item \kotlincode{6.666\ldots}
            \item \kotlincode{6}
            \item \kotlincode{7}
            \item \kotlincode{20 / 3}
        \end{enumerate}
        \quizpause
    \end{frame}

    \begin{frame}{Answer --- Quiz 2}
        \answerbox{\textbf{2.} \kotlincode{6}}
        \vspace{0.4em}
        Both operands are \kotlincode{Int}, so the result is truncated integer division.
        Remainder would be \kotlincode{20 \% 3} $= 2$.
    \end{frame}

%======================================================================
    \section{Mutable variables}
%======================================================================

    \begin{frame}{\kotlincode{val} vs \kotlincode{var}}
        \begin{columns}[T]
            \begin{column}{0.48\textwidth}
                \textbf{\kotlincode{val} (immutable binding)}
                \begin{itemize}
                    \item Assign once
                    \item Cannot reassign
                    \item Prefer by default
                \end{itemize}
            \end{column}
            \begin{column}{0.48\textwidth}
                \textbf{\kotlincode{var} (mutable binding)}
                \begin{itemize}
                    \item Can reassign
                    \item Same name, new value
                    \item Use when the value must change
                \end{itemize}
            \end{column}
        \end{columns}
        \vspace{0.8em}
        ``Immutable binding'' means the \emph{variable} cannot be pointed at a new value.
        (Contents of a mutable list can still change --- later.)
    \end{frame}

    \begin{frame}[fragile]{Updating with \kotlincode{var}}
        \begin{lstlisting}
fun main() {
    var score = 0
    score = score + 10
    score += 5          // same as score = score + 5
    println(score)      // 15

    var lives = 3
    lives -= 1
    println(lives)      // 2
}
        \end{lstlisting}
        \vspace{0.3em}
        Compound assignments: \kotlincode{+=}, \kotlincode{-=}, \kotlincode{*=}, \kotlincode{/=}, \kotlincode{\%=}
    \end{frame}

    \begin{frame}[fragile]{When \kotlincode{val} is enough}
        \begin{lstlisting}
fun main() {
    val a = readln().toInt()
    val b = readln().toInt()
    val sum = a + b       // no need for var
    println(sum)
}
        \end{lstlisting}
        \vspace{0.3em}
        \commonmistake{Using \kotlincode{var} everywhere ``just in case''.
        Start with \kotlincode{val}; switch to \kotlincode{var} only when reassignment is required.}
    \end{frame}

    %--- Quiz checkpoint 3 ---
    \begin{frame}{Quiz 3 --- Mutable variables}
        \quizbox{Which snippet compiles?}
        \begin{enumerate}
            \item \kotlincode{val x = 1; x = 2}
            \item \kotlincode{var x = 1; x = 2}
            \item \kotlincode{val x; x = 1}
            \item \kotlincode{var x = 1; x = "hi"}
        \end{enumerate}
        \quizpause
    \end{frame}

    \begin{frame}{Answer --- Quiz 3}
        \answerbox{\textbf{2.} \kotlincode{var x = 1; x = 2}}
        \vspace{0.4em}
        (1) cannot reassign \kotlincode{val}.
        (3) needs a type or initializer.
        (4) type is fixed as \kotlincode{Int}; cannot assign a \kotlincode{String}.
    \end{frame}

%======================================================================
    \section{Floating-point numbers}
%======================================================================

    \begin{frame}{Floating-point types}
        For numbers with a fractional part, use floating-point types:
        \vspace{0.5em}
        \begin{tabular}{lll}
            \toprule
            Type & Size & Typical use \\
            \midrule
            \kotlincode{Double} & 64-bit & default decimal type \\
            \kotlincode{Float} & 32-bit & less precision; rarer in intro code \\
            \bottomrule
        \end{tabular}
        \vspace{0.6em}
        Literals: \kotlincode{3.14} is \kotlincode{Double}; \kotlincode{3.14f} is \kotlincode{Float}.
    \end{frame}

    \begin{frame}[fragile]{Working with \kotlincode{Double}}
        \begin{lstlisting}
fun main() {
    val price = 19.99
    val taxRate = 0.08
    val total = price * (1 + taxRate)
    println("Total: $total")

    print("Enter a decimal: ")
    val x = readln().toDouble()
    println("Half is ${x / 2}")
}
        \end{lstlisting}
    \end{frame}

    \begin{frame}{Precision warning}
        Floating-point arithmetic is approximate in binary.
        \vspace{0.5em}
        \begin{itemize}
            \item \kotlincode{0.1 + 0.2} may not be exactly \kotlincode{0.3}
            \item Fine for science demos and many apps; careful with money
            \item For currency, later courses use integer cents or decimal libraries
        \end{itemize}
        \vspace{0.4em}
        \trythis{Print \kotlincode{0.1 + 0.2} and discuss what you see.}
    \end{frame}

    %--- Quiz checkpoint 4 ---
    \begin{frame}{Quiz 4 --- Floating-point}
        \quizbox{How do you read a decimal number from the console into a \kotlincode{Double}?}
        \begin{enumerate}
            \item \kotlincode{val x = readln()}
            \item \kotlincode{val x = readln().toInt()}
            \item \kotlincode{val x = readln().toDouble()}
            \item \kotlincode{val x = Double(readln())}
        \end{enumerate}
        \quizpause
    \end{frame}

    \begin{frame}{Answer --- Quiz 4}
        \answerbox{\textbf{3.} \kotlincode{val x = readln().toDouble()}}
        \vspace{0.4em}
        Parse the input string with \kotlincode{toDouble()}.
        \kotlincode{toInt()} would reject values like \kotlincode{"3.5"}.
    \end{frame}

%======================================================================
    \section{Long, Short, Byte}
%======================================================================

    \begin{frame}{Other integer types}
        \begin{tabular}{llll}
            \toprule
            Type & Bits & Approx.\ range & Literal tip \\
            \midrule
            \kotlincode{Byte} & 8 & $-128$ \ldots\ $127$ & convert explicitly \\
            \kotlincode{Short} & 16 & $\approx \pm 32\,000$ & convert explicitly \\
            \kotlincode{Int} & 32 & $\approx \pm 2\cdot 10^9$ & default \\
            \kotlincode{Long} & 64 & very large & suffix \kotlincode{L} \\
            \bottomrule
        \end{tabular}
        \vspace{0.6em}
        For most intro programs, \kotlincode{Int} and \kotlincode{Long} are enough.
    \end{frame}

    \begin{frame}[fragile]{\kotlincode{Long} in practice}
        \begin{lstlisting}
fun main() {
    val population = 8_000_000_000L  // underscores OK
    val big = 2_000_000_000
    // big * 3 may overflow Int; use Long:
    val safer = big.toLong() * 3
    println(population)
    println(safer)
}
        \end{lstlisting}
        \vspace{0.3em}
        \kotlincode{toLong()}, \kotlincode{toShort()}, \kotlincode{toByte()} convert when needed.
        Converting out of range can wrap or throw depending on API --- stay in range.
    \end{frame}

    \begin{frame}[fragile]{\kotlincode{Byte} and \kotlincode{Short}}
        \begin{lstlisting}
fun main() {
    val small: Byte = 100
    val medium: Short = 30_000
    // Arithmetic often promotes to Int:
    val sum = small + medium   // Int
    println(sum)
}
        \end{lstlisting}
        \vspace{0.3em}
        You rarely declare \kotlincode{Byte}/\kotlincode{Short} unless matching a file format,
        network protocol, or memory-tight structure.
    \end{frame}

    %--- Quiz checkpoint 5 ---
    \begin{frame}{Quiz 5 --- Integer sizes}
        \quizbox{Which literal is a \kotlincode{Long}?}
        \begin{enumerate}
            \item \kotlincode{42}
            \item \kotlincode{42.0}
            \item \kotlincode{42L}
            \item \kotlincode{"42"}
        \end{enumerate}
        \quizpause
    \end{frame}

    \begin{frame}{Answer --- Quiz 5}
        \answerbox{\textbf{3.} \kotlincode{42L}}
        \vspace{0.4em}
        The \kotlincode{L} suffix marks a \kotlincode{Long} literal.
        \kotlincode{42} is \kotlincode{Int}; \kotlincode{42.0} is \kotlincode{Double}.
    \end{frame}

%======================================================================
    \section{Wrap-up}
%======================================================================

    \begin{frame}{Summary}
        \begin{itemize}
            \item \kotlincode{Int} --- default whole-number type
            \item Integer \kotlincode{/} truncates; use \kotlincode{\%} for remainder
            \item \kotlincode{var} when you must reassign; prefer \kotlincode{val} otherwise
            \item \kotlincode{Double} for decimals; watch floating-point precision
            \item \kotlincode{Long} / \kotlincode{Short} / \kotlincode{Byte} for other integer widths
        \end{itemize}
    \end{frame}

    \begin{frame}{Looking ahead}
        Next: conditional operators (\kotlincode{if}/\kotlincode{else}), combining conditions,
        and \kotlincode{if}/\kotlincode{when} as expressions.
    \end{frame}

\end{document}
